\documentclass[11pt,a4paper]{article}
% preamble.tex — estilo común de todos los apuntes Froth.
% Cada apunte hace % preamble.tex — estilo común de todos los apuntes Froth.
% Cada apunte hace % preamble.tex — estilo común de todos los apuntes Froth.
% Cada apunte hace \input{preamble} justo después de \documentclass.
\usepackage[margin=2.4cm]{geometry}
\usepackage{xcolor}
\definecolor{litblue}{RGB}{31,79,121}
\definecolor{litgreen}{RGB}{27,94,32}
\definecolor{litorange}{RGB}{230,81,0}
\usepackage{parskip}
\usepackage{titlesec}
\titleformat{\section}{\Large\bfseries\color{litblue}}{\thesection}{0.6em}{}
\titleformat{\subsection}{\large\bfseries\color{litblue}}{\thesubsection}{0.6em}{}
\usepackage{tcolorbox}
\usepackage{fancyhdr}
\pagestyle{fancy}
\fancyhf{}
\fancyhead[L]{\small\color{litblue}Froth — Apuntes de ML}
\fancyhead[R]{\small\thepage}
\renewcommand{\headrulewidth}{0.4pt}
\usepackage[colorlinks=true,linkcolor=litblue,urlcolor=litblue]{hyperref}

% Cabecera de cada apunte: \cabecera{numero}{titulo}
\newcommand{\cabecera}[2]{%
  \begin{tcolorbox}[colback=litblue,coltext=white,colframe=litblue,arc=2mm]
    {\Large\bfseries Apunte #1 — #2}\\[3pt]
    {\small Proyecto Froth · aprender Machine Learning construyendo}
  \end{tcolorbox}\vspace{0.8em}}

% Cajas temáticas
\newtcolorbox{concepto}[1]{colback=blue!4,colframe=litblue,title={Concepto: #1},arc=1.5mm}
\newtcolorbox{practica}{colback=green!5,colframe=litgreen,title={Pruébalo tú (teclado en mano)},arc=1.5mm}
\newtcolorbox{ojo}{colback=orange!8,colframe=litorange,title={Ojo / error típico},arc=1.5mm}
\newtcolorbox{resumen}{colback=gray!8,colframe=gray!60,title={En una frase},arc=1.5mm}

\newcommand{\codigo}[1]{\texttt{#1}}
 justo después de \documentclass.
\usepackage[margin=2.4cm]{geometry}
\usepackage{xcolor}
\definecolor{litblue}{RGB}{31,79,121}
\definecolor{litgreen}{RGB}{27,94,32}
\definecolor{litorange}{RGB}{230,81,0}
\usepackage{parskip}
\usepackage{titlesec}
\titleformat{\section}{\Large\bfseries\color{litblue}}{\thesection}{0.6em}{}
\titleformat{\subsection}{\large\bfseries\color{litblue}}{\thesubsection}{0.6em}{}
\usepackage{tcolorbox}
\usepackage{fancyhdr}
\pagestyle{fancy}
\fancyhf{}
\fancyhead[L]{\small\color{litblue}Froth — Apuntes de ML}
\fancyhead[R]{\small\thepage}
\renewcommand{\headrulewidth}{0.4pt}
\usepackage[colorlinks=true,linkcolor=litblue,urlcolor=litblue]{hyperref}

% Cabecera de cada apunte: \cabecera{numero}{titulo}
\newcommand{\cabecera}[2]{%
  \begin{tcolorbox}[colback=litblue,coltext=white,colframe=litblue,arc=2mm]
    {\Large\bfseries Apunte #1 — #2}\\[3pt]
    {\small Proyecto Froth · aprender Machine Learning construyendo}
  \end{tcolorbox}\vspace{0.8em}}

% Cajas temáticas
\newtcolorbox{concepto}[1]{colback=blue!4,colframe=litblue,title={Concepto: #1},arc=1.5mm}
\newtcolorbox{practica}{colback=green!5,colframe=litgreen,title={Pruébalo tú (teclado en mano)},arc=1.5mm}
\newtcolorbox{ojo}{colback=orange!8,colframe=litorange,title={Ojo / error típico},arc=1.5mm}
\newtcolorbox{resumen}{colback=gray!8,colframe=gray!60,title={En una frase},arc=1.5mm}

\newcommand{\codigo}[1]{\texttt{#1}}
 justo después de \documentclass.
\usepackage[margin=2.4cm]{geometry}
\usepackage{xcolor}
\definecolor{litblue}{RGB}{31,79,121}
\definecolor{litgreen}{RGB}{27,94,32}
\definecolor{litorange}{RGB}{230,81,0}
\usepackage{parskip}
\usepackage{titlesec}
\titleformat{\section}{\Large\bfseries\color{litblue}}{\thesection}{0.6em}{}
\titleformat{\subsection}{\large\bfseries\color{litblue}}{\thesubsection}{0.6em}{}
\usepackage{tcolorbox}
\usepackage{fancyhdr}
\pagestyle{fancy}
\fancyhf{}
\fancyhead[L]{\small\color{litblue}Froth — Apuntes de ML}
\fancyhead[R]{\small\thepage}
\renewcommand{\headrulewidth}{0.4pt}
\usepackage[colorlinks=true,linkcolor=litblue,urlcolor=litblue]{hyperref}

% Cabecera de cada apunte: \cabecera{numero}{titulo}
\newcommand{\cabecera}[2]{%
  \begin{tcolorbox}[colback=litblue,coltext=white,colframe=litblue,arc=2mm]
    {\Large\bfseries Apunte #1 — #2}\\[3pt]
    {\small Proyecto Froth · aprender Machine Learning construyendo}
  \end{tcolorbox}\vspace{0.8em}}

% Cajas temáticas
\newtcolorbox{concepto}[1]{colback=blue!4,colframe=litblue,title={Concepto: #1},arc=1.5mm}
\newtcolorbox{practica}{colback=green!5,colframe=litgreen,title={Pruébalo tú (teclado en mano)},arc=1.5mm}
\newtcolorbox{ojo}{colback=orange!8,colframe=litorange,title={Ojo / error típico},arc=1.5mm}
\newtcolorbox{resumen}{colback=gray!8,colframe=gray!60,title={En una frase},arc=1.5mm}

\newcommand{\codigo}[1]{\texttt{#1}}

\begin{document}
\cabecera{02}{APIs y JSON: cómo se piden datos a un servidor}

\section{La idea en 30 segundos}

Tu PC (el \emph{cliente}) le manda una petición por internet a otra computadora
(el \emph{servidor}) y esta responde con datos. Eso es todo. Una \textbf{API}
(\emph{Application Programming Interface}) es simplemente la ventanilla oficial que un
servicio ofrece para que \emph{programas} (no personas con navegador) le pidan cosas.

En Froth usamos la API de \textbf{OpenAlex}, un catálogo gratuito con $\sim$250 millones
de trabajos académicos.

\begin{concepto}{petición HTTP GET}
Una petición se escribe como una URL con \textbf{parámetros} después del signo \codigo{?}:
\begin{center}
\small\codigo{https://api.openalex.org/works?search=lepidolite flotation\&per\_page=25}
\end{center}
Se lee: ``del catálogo \codigo{works}, búscame \codigo{lepidolite flotation}, dame 25 por
página''. En Python, la librería \codigo{requests} construye y envía esto por ti:
\begin{center}
\codigo{r = requests.get(url, params=\{...\})}
\end{center}
\end{concepto}

\begin{concepto}{JSON}
El servidor responde en \textbf{JSON}: un formato de texto con la misma pinta que los
diccionarios y listas de Python (\codigo{\{"title": "...", "year": 2021\}}).
Por eso \codigo{r.json()} lo convierte directamente en objetos de Python con los que ya
puedes trabajar. JSON es \emph{el} idioma de intercambio de datos en internet.
\end{concepto}

\section{Detalles de ciudadano educado}

Las APIs públicas son un recurso compartido. Tres costumbres que nuestro \codigo{pull.py}
ya practica:

\begin{enumerate}
  \item \textbf{Identifícate} (\emph{polite pool}): mandamos nuestro email en el parámetro
        \codigo{mailto}. OpenAlex da mejor servicio a quien da la cara.
  \item \textbf{No bombardees}: entre búsqueda y búsqueda esperamos 0.3 s
        (\codigo{time.sleep(0.3)}). Miles de peticiones por segundo tumban servicios.
  \item \textbf{Pide solo lo que necesitas}: el parámetro \codigo{select} lista los campos
        que queremos (título, año, autores...). Menos datos $=$ más rápido para todos.
\end{enumerate}

\begin{concepto}{el truco del abstract invertido}
OpenAlex no guarda el resumen como texto corrido sino como un \emph{índice invertido}:
un diccionario \codigo{\{palabra: [posiciones donde aparece]\}}. Es una curiosidad real de
esta API (lo hacen por temas legales y de eficiencia). Nuestra función
\codigo{\_reconstruct\_abstract()} lo vuelve a montar: coloca cada palabra en su posición
y las une. Primer ejemplo del proyecto de que \textbf{los datos reales vienen ``sucios''}
y hay que trabajarlos.
\end{concepto}

\begin{practica}
Pega esta URL en tu navegador (¡una API también responde ahí!) y mira el JSON crudo:
\begin{center}
\small\codigo{https://api.openalex.org/works?search=lepidolite\%20flotation\&per\_page=2}
\end{center}
Busca los campos \codigo{title}, \codigo{publication\_year} y
\codigo{abstract\_inverted\_index} en la respuesta.
\end{practica}

\begin{resumen}
Una API es una ventanilla para programas: mandas una URL con parámetros, te responden
JSON, y \codigo{requests} + \codigo{r.json()} lo convierten en diccionarios de Python.
\end{resumen}

\end{document}
